% Copyright (c) 2026 Eugene Joseph M. Ragasa.
% SPDX-License-Identifier: Apache-2.0
% Author: Eugene Joseph M. Ragasa
% Formal proof completed: 2026-10-20
\documentclass[11pt]{article}
\usepackage[T1]{fontenc}
\usepackage{amsmath,amssymb,booktabs,tabularx}
\usepackage[margin=1in]{geometry}
\usepackage{xurl}
\usepackage[hidelinks]{hyperref}
\setlength{\emergencystretch}{3em}
\usepackage{listings}
\lstset{basicstyle=\ttfamily\small,breaklines=true,columns=fullflexible,
  literate={ℝ}{{$\mathbb R$}}1 {ℂ}{{$\mathbb C$}}1 {¬}{{$\neg$}}1
           {∃}{{$\exists$}}1 {≠}{{$\ne$}}1}
\newcommand{\StatementID}{BG-PROJECTION-THM-002}
\newcommand{\LeanDeclaration}{ProjectKoios.BandGauge.Projection.twoBandLaurentRadicand\_not\_sq}
\title{Project Koios Formal Proof Dossier\\\large Laurent Square-Root Obstruction for the Two-Band Radicand}
\author{Eugene Joseph M. Ragasa}
\date{Formal proof completed: 2026-10-20}
\begin{document}
\sloppy
\maketitle

\section{Identity and status}
\begin{tabularx}{\textwidth}{@{}lX@{}}\toprule
Statement ID & \path{BG-PROJECTION-THM-002} \\
Status & proved; no admissions \\
Lean declaration & \path{ProjectKoios.BandGauge.Projection.twoBandLaurentRadicand_not_sq} \\
Owning module & \path{ProjectKoios/BandGauge/Projection/LaurentObstruction.lean} \\
Author and license & Eugene Joseph M. Ragasa; \texttt{Apache-2.0} \\
Formal proof completed & 2026-10-20 \\
\bottomrule\end{tabularx}

This status applies only to the exact Laurent-polynomial non-square theorem below. The
complete publication proposition is proved separately by \path{BG-PROJECTION-THM-008}.

\section{Exact mathematical statement}
For $\Delta,t\in\mathbb R$ with $\Delta t\ne0$, define
\[
 r(z)=\Delta^2+2t^2+t^2z+t^2z^{-1}\in\mathbb C[z,z^{-1}].
\]
Then no $f\in\mathbb C[z,z^{-1}]$ satisfies $f^2=r$:
\[
 \neg\,\exists f\in\mathbb C[z,z^{-1}],\qquad f(z)^2=r(z).
\]

\section{Source provenance}
\begin{tabularx}{\textwidth}{@{}lX@{}}\toprule
Owner repository & \texttt{ksdft2effmass} \\
Navigation branch & \texttt{work/paper-3-higher-rank-extension} \\
Immutable commit & \path{99a46cb59d0c6f27ada7d709f194996ef9619d9f} \\
Statement path and locator & \path{docs/publications/papers/Ragasa2026b/v2/manuscript.tex:316--340} \\
Proof path and locator & \path{docs/publications/papers/Ragasa2026b/v2/appendix_proofs.tex:11--24} \\
Statement SHA-256 & \path{d8d5951d8d7ed75355b5d94f37fcbddc70c1466897a5770dbe9f877e60016daf} \\
Proof SHA-256 & \path{614fad3426a941a9a1c5ad4c73ca917b55007b525a1a3a76970d5d7df5af7d19} \\
Dirty-state limitation & cited bytes are the clean immutable commit recorded by the owner-source audit \\
\bottomrule\end{tabularx}

The declaration \emph{formalizes} the Laurent-polynomial contradiction in the appendix.
It does not formalize the preceding implication from a finite Fourier function on the
unit circle to a Laurent-polynomial identity.

\section{Assumptions, conventions, and statement delta}
\begin{tabularx}{\textwidth}{@{}XXX@{}}\toprule
Source assumption or convention & Lean encoding & Delta or limitation \\
\midrule
$\Delta t\ne0$ & exact real premise \texttt{gap * hopping != 0} & preserved exactly \\
Finite Fourier expression in $z=e^{ik}$ & finite-support
\texttt{LaurentPolynomial $\mathbb C$} indexed by $\mathbb Z$ & this theorem begins only
after the Laurent identity has been obtained \\
Right-hand support $\{-1,0,1\}$ & exact three-term Laurent definition & only the
positive extreme is used by the Lean contradiction \\
Complex coefficients for a hypothetical square root &
\texttt{f : LaurentPolynomial $\mathbb C$} & matches the stronger complex-coefficient
possibility, hence also excludes real coefficients \\
Source proof mentions smallest and largest exponents & Lean proof uses the largest
exponent and the nonzero coefficient at exponent $1$ & same conclusion and source
assumption; no omitted case \\
\bottomrule\end{tabularx}

\section{Statement-specific prior art}
\begin{tabularx}{\textwidth}{@{}lXXX@{}}\toprule
Relation & Canonical identity & Locator & Role and caveat \\
\midrule
formalizes & owner source above; canonical literature identity unresolved &
\path{appendix_proofs.tex:11--24} & supplies the stated extreme-support argument \\
background-only & unresolved & unverified & finite-support convolution and parity are
standard algebra; no novelty or priority claim is made \\
\bottomrule\end{tabularx}

The References owner has not supplied an accepted external prior-art identity for this
specific elementary obstruction. A bounded search or absence of a citation would not
establish originality.

\section{Motivation and contribution boundary}
The theorem isolates the exact algebraic obstruction behind the source counterexample:
a Laurent square has an even largest exponent, whereas the radicand has a nonzero term at
exponent $1$ and no terms above it. The contribution here is a traceable, kernel-checked
encoding of that argument, not a claim that the algebraic idea is new.

\section{Proof schematic}
\[
\boxed{f^2=r,\ \Delta t\ne0}
\Longrightarrow
\boxed{[z^1]r=t^2\ne0\ \Longrightarrow\ 1\in\operatorname{supp}(f^2)}
\Longrightarrow
\boxed{1\le2n,\quad n=\max\operatorname{supp}(f)}
\]
\[
\boxed{[z^{2n}]f^2=([z^n]f)^2\ne0}
\Longrightarrow
\boxed{2n\le1}
\Longrightarrow
\boxed{2n=1}
\Longrightarrow\bot.
\]
The second boxed implication uses that every coefficient of $r$ above exponent $1$ is
zero. The last contradiction is integer parity.

\section{Human-readable mathematical argument}
Assume that $f^2=r$. Since $\Delta t\ne0$, in particular $t\ne0$, and the coefficient
of $z^1$ in $r$ is $t^2\ne0$. Thus $f$ is nonzero. Let $n$ be the largest exponent with
nonzero coefficient in $f$.

Because the coefficient of $z^1$ in $f^2$ is nonzero, the convolution-support inclusion
provides exponents $a,b$ in the support of $f$ with $a+b=1$. Both are at most $n$, so
$1\le2n$.

At exponent $2n$, the only support pair that can contribute to the convolution is
$(n,n)$: if $x<n$ then $2n-x>n$ and its coefficient vanishes. Hence
$[z^{2n}]f^2=([z^n]f)^2\ne0$. But $r$ has no nonzero coefficient above exponent $1$,
so $2n\le1$. Combining the inequalities gives $2n=1$, impossible for an integer $n$.
Therefore no such Laurent polynomial exists.

\section{Lean realization}
\begin{lstlisting}
theorem twoBandLaurentRadicand_not_sq
    (gap hopping : ℝ) (gap_mul_hopping_ne : gap * hopping ≠ 0) :
    ¬ ∃ f : LaurentPolynomial ℂ,
      f ^ 2 = twoBandLaurentRadicand gap hopping
\end{lstlisting}
The private implementation lemma computes the extreme convolution coefficient. The main
proof uses \texttt{support\_coeff\_mul\_subset}, finite-support maximum lemmas, and exact
integer arithmetic.

\noindent
\begin{tabularx}{\textwidth}{@{}lX@{}}\toprule
Lean toolchain & Lean 4.35.0 \\
Mathlib revision & \texttt{56a005855ea17e0b0958f73ec801ae140b0a1679} \\
Additional dependencies & none; Physlib not adopted \\
Principal imports & Laurent polynomials and add-monoid-algebra degree/support infrastructure \\
Validation & see \path{docs/trusted-base-audit.md} for historical exact-source and named-audit artifacts \\
Promotion & current reconciliation tree still requires promotion gates \\
\bottomrule\end{tabularx}

\section{Trusted base and axiom report}
\begin{lstlisting}
#print axioms ProjectKoios.BandGauge.Projection.twoBandLaurentRadicand_not_sq
-- [propext, Classical.choice, Quot.sound]
\end{lstlisting}
There are no repository-owned axioms or admissions in the declaration.

\section{AI, automation, and comprehension record}
The theorem selection, proof search, Lean implementation, and dossier prose were
AI-assisted. Lean elaboration and kernel checking validated the exact declaration.
Simplification, \texttt{norm\_num}, and \texttt{omega} produced proof terms checked by
Lean. The maintainer selected the source passage and reviewed the support/parity argument
and the stated boundaries. No independent expert mathematical review has occurred.

\section{Scientific and evidence boundaries}
This exact theorem has no numerical premise and consumes no runtime or scientific
evidence. By itself it does not show that a finite Fourier function equal to the positive
square root yields the displayed Laurent identity or identify $P_-HP_-$. Those bridges and
the final matrix-valued conclusion are proved separately by
\path{BG-PROJECTION-THM-003--008}. It says nothing about a material Hamiltonian.

\section{Unresolved questions}
\begin{itemize}
\item Resolve statement-specific external prior-art identity through References.
\item Obtain independent expert review if this proof package becomes publication-facing.
\item Quantitative Fourier decay and truncation error remain outside the exact
counterexample.
\end{itemize}

\section{Review and revision history}
\begin{tabularx}{\textwidth}{@{}lllX@{}}\toprule
Date & Revision & Reviewer or process & Outcome and limitations \\
\midrule
2026-10-20 & initial & AI-assisted implementation; Lean kernel check & Laurent
square-root obstruction proved under the source assumption; independent expert review
remains open \\
2026-10-20 & companion integration & Lean kernel check & unit-circle, projector, and final
matrix bridges completed separately through \path{BG-PROJECTION-THM-008} \\
\bottomrule\end{tabularx}
\end{document}
